\documentclass[12pt]{article}
\usepackage[utf8]{inputenc}
\usepackage[T1]{fontenc}
\usepackage[margin=1in]{geometry}
\usepackage{setspace}
\usepackage{xcolor}
\usepackage[colorlinks=true, urlcolor=blue!60!black, linkcolor=blue!60!black]{hyperref}
\setlength{\parindent}{0em}
\setlength{\parskip}{0.6em}
\renewcommand{\baselinestretch}{1.0}
\begin{document}
\begin{center}
{\LARGE \textit{FDA v. Alliance for Hippocratic Medicine}, 602 U.S. 367 (2024)\par}
\vspace{0.5em}
{\large Case Brief --- Two-Page Edition\par}
\end{center}

\section{Memory Jogger}
Doctors and groups who neither prescribe nor use mifepristone lacked standing to challenge the FDA's relaxed rules; an organization ``cannot spend its way into standing'' (\href{https://www.supremecourt.gov/opinions/23pdf/602us1r35_bqmd.pdf}{602 U.S. 367}).

\section{Detailed Case Facts}
In 2016 and 2021 the FDA eased restrictions on mifepristone, extending its use to ten weeks, allowing non-physician prescribers, and permitting mail dispensing. Pro-life medical associations and doctors sued under the APA, though none prescribed or used the drug and the FDA required nothing of them.

\section{Procedural History}
The district court effectively removed mifepristone from the market; the Supreme Court stayed that order. The Fifth Circuit found standing and upheld relief against the 2016 and 2021 actions.

\section{Judicial Votes}
Reversed, 9--0. Kavanaugh wrote for the Court; Thomas concurred.

\section{Holding}
Unregulated plaintiffs lack standing to challenge the FDA's deregulation of a drug they do not prescribe or use. Conscience laws break the causal chain, downstream costs are too attenuated, resource diversion does not create organizational standing, and the absence of other plaintiffs is irrelevant.

\section{Analysis of Opinions}
\textbf{Majority:} When plaintiffs challenge the regulation of others, causation is especially hard to establish. Federal conscience laws protect doctors, and none had ever been forced to act against conscience. Recognizing doctor standing would let anyone who treats injuries challenge any safety deregulation. \textit{Havens Realty} involved direct interference with core business and is not to be extended. \textbf{Thomas:} associational standing, which lets an uninjured group sue for one member, deserves reconsideration.

\section{Examples: Future Applications}
\textbf{Same-Side:} advocacy groups that divert resources to fight a voter-roll law lack standing. \textbf{Opposite-Side:} a rival drug maker challenging an unlawful competitor approval has standing. \textbf{Fence-Sitter:} a legal-aid nonprofit whose attorneys are barred from detained clients.

\section{Critique}
Conservatives praise a disciplined Article III ruling. Critics note that the merits remain open to state plaintiffs, contrast the Court's looser standing in \textit{Biden v. Nebraska}, and warn that narrowing \textit{Havens} will hobble civil-rights nonprofits.

\section{Key Quotations}
\begin{quote}
``But an organization that has not suffered a concrete injury caused by a defendant's action cannot spend its way into standing simply by expending money to gather information and advocate against the defendant's action. An organization cannot manufacture its own standing in that way.''
\end{quote}
\noindent --- Justice Kavanaugh, for the Court, 602 U.S. at 394.

\end{document}
